\documentclass[11pt,a4paper]{article}
\usepackage[T1]{fontenc}
\usepackage[utf8]{inputenc}

\usepackage[margin=1in,left=0.75in,right=0.7in,top=0.5in,bottom=0.45in]{geometry}
\usepackage{xcolor}
\usepackage{tikz}
\usepackage{parskip}
\usepackage{enumitem}
\usepackage{needspace}
\usepackage{eso-pic}
\usepackage{listings}
\usepackage{tabularx}
\usepackage{booktabs}
\usepackage{graphicx}
\usepackage[colorlinks=true,urlcolor=teal,linkcolor=teal]{hyperref}

\definecolor{ink}{HTML}{111827}
\definecolor{teal}{HTML}{0F766E}
\definecolor{sand}{HTML}{FBFAF6}
\definecolor{rule}{HTML}{D9D6CC}
\definecolor{muted}{HTML}{6B7280}
\definecolor{warn}{HTML}{B45309}
\pagecolor{sand}\color{ink}
\pagestyle{empty}

\AddToShipoutPictureBG{\begin{tikzpicture}[remember picture,overlay]
  \fill[teal] (current page.north west) rectangle ([xshift=0.35in]current page.south west);
\end{tikzpicture}}

\newcommand{\kicker}[1]{{\sffamily\bfseries\small\color{teal}\MakeUppercase{#1}}}
\newcommand{\lbl}[1]{{\sffamily\bfseries\footnotesize\color{muted}\MakeUppercase{#1}}}
\newcommand{\sect}[1]{\needspace{5\baselineskip}\vspace{0.5em}\par\kicker{#1}\par\vspace{-0.5em}{\color{rule}\rule{\linewidth}{0.5pt}}\par\vspace{-0.2em}}
\newcommand{\todo}[1]{{\color{warn}\bfseries[TODO: #1]}}
\setlength{\parskip}{0.3em}
\setlist[itemize]{itemsep=1pt,topsep=1pt,leftmargin=1.2em}

\lstset{
  basicstyle=\ttfamily\footnotesize\color{ink},
  breaklines=true,
  breakatwhitespace=true,
  frame=single,
  rulecolor=\color{rule},
  backgroundcolor=\color{white},
  framesep=5pt,
  xleftmargin=3pt,
  xrightmargin=3pt,
  columns=fullflexible,
  keepspaces=true,
  aboveskip=4pt,
  belowskip=4pt
}

\newcommand{\notes}[4]{%
\lbl{Notes}\par
\begin{itemize}
  \item \textbf{What changed in the prompt:} #1
  \item \textbf{What improved in the output:} #2
  \item \textbf{What still failed:} #3
  \item \textbf{What I would try next:} #4
\end{itemize}}
\newcommand{\fig}[2]{%
\begin{figure}[htbp]
\centering
\includegraphics[width=\linewidth]{figures/#1}
\caption{#2}
\end{figure}}

\begin{document}

\kicker{ShadowFox data analyst internship}
\vspace{0.2em}
\begin{flushleft}
{\LARGE\bfseries Azizul Abedin Azmi}\par
{\small\color{muted} \href{mailto:azmi00121@gmail.com}{azmi00121@gmail.com} \quad|\quad \href{https://github.com/azizulabedinazmi}{GitHub} \quad|\quad \href{https://azizulabedinazmi.is-a.dev/}{Portfolio}}
\end{flushleft}
\begin{flushright}
{\sffamily\bfseries\fontsize{25}{28}\selectfont Beginner\\
Superstore analysis}\\[0.4em]
{\color{muted}\sffamily Rebuild, validation and findings report}\\
{\color{muted}\sffamily 11 October 2026}
\end{flushright}

\sect{Executive summary}
The Level 1 project cleans the Sample Superstore export and builds a live Excel
dashboard from 9,993 order lines, 5,009 orders, 793 customers and 1,862 product
IDs. The complete pipeline now runs in the documented order and passes a
regression validator. The resulting workbook contains KPI definitions, monthly
trend and seasonality, product and discount analysis, a dashboard, and an
auditable cleaning log.

\begin{tabularx}{\linewidth}{@{}lXr@{}}
\toprule
\textbf{Metric} & \textbf{Definition} & \textbf{Result}\\
\midrule
Sales & Sum of cleaned line sales & \$2,296,919.49\\
Profit & Sum of cleaned line profit & \$286,409.08\\
Profit margin & Profit divided by sales & 12.5\%\\
Orders & Distinct order IDs & 5,009\\
Return rate & Returned orders divided by orders & 5.9\%\\
\bottomrule
\end{tabularx}

\fig{dashboard_overview.png}{Dashboard-style overview reconstructed from the cleaned data. The seven KPI cards and the monthly sales/profit curve correspond to the main Excel Dashboard sheet.}

\sect{Workbook dashboard and visual evidence}
The Excel workbook is not only a table of results: it is a formula-driven
dashboard. The report therefore includes the same visual story in a portable
PDF form. Values in these figures are recalculated from
\texttt{data/clean/superstore\_clean.csv}; the workbook remains the interactive
version with formulas, charts and filters.

\fig{monthly_trend.png}{Monthly sales curve and three-month moving average. The curve shows recurring year-end peaks and a clear increase in the later years.}

\fig{category_region.png}{Sales and profit by category and region. Furniture is much less profitable than Office Supplies and Technology, while the West and East carry the largest sales totals.}

\sect{What was fixed}
\begin{itemize}
  \item The cleaner now stops with an explicit error when any integrity check
  fails, instead of writing ``FAILED'' to the log and continuing.
  \item ``First Line of Order'' now uses a running \texttt{COUNTIF}, so order
  counts remain correct even if an input file is not grouped by order.
  \item A new \texttt{validate\_project.py} script checks row counts, totals,
  data constraints, required sheets, formulas and dashboard structure.
  \item The README source path was corrected, the screenshot placeholder was
  replaced with usable workbook instructions, and recommendations were added.
\end{itemize}

\sect{Reproducible workflow}
\begin{lstlisting}
cd "d:\ShadowFox Intership\shadowfox-data-analyst\01-beginner"
python scripts\clean_superstore.py
python scripts\build_workbook.py
python scripts\add_analysis.py
python scripts\add_dashboard.py
python scripts\validate_project.py
\end{lstlisting}
The raw CSV is preserved. The first script splits the combined Orders, People
and Returns export, converts types, fills the 11 missing Burlington ZIP codes,
removes one duplicate line and records every action in the cleaning log. The
three workbook-building scripts are intentionally ordered because later scripts
add sheets and formulas to the workbook produced by earlier scripts.

\sect{Business findings}
\begin{enumerate}[leftmargin=1.5em,itemsep=2pt]
  \item Deep discounting is the clearest margin warning. Lines above 20\%
  discount are loss-making in aggregate, so discount approval should be
  reviewed by sub-category rather than applied broadly.
  \item Furniture is the weakest category. Tables and Bookcases are the main
  sub-category loss areas; Machines also combines substantial sales with a low
  margin and a high average discount.
  \item Demand is seasonal. September, November and December contribute about
  43\% of sales, while January and February contribute about 7\%. The workbook
  compares calendar months across all four years.
  \item Product concentration matters. A small number of products account for
  very large gains or losses, making product-level review more useful than
  product-ID-only grouping because 32 IDs map to multiple names.
  \item Returns are material but not dominant: 5.9\% of orders and 7.9\% of
  sales are associated with returned orders. Headline figures remain gross
  because the source does not state whether profit was reversed.
\end{enumerate}

\sect{Category, product and margin detail}
\fig{subcategory_performance.png}{Sub-category performance. Red bars identify negative profit or negative margin; this makes Tables and Bookcases immediately visible as priority areas.}

\fig{product_rankings.png}{Top and bottom ten products by aggregate profit. Product names are used rather than Product IDs because the source reuses 32 IDs for different names.}

\fig{discount_analysis.png}{Discount analysis. The left chart groups lines into the workbook's four discount bands; the right chart shows the relationship at each exact discount rate.}

\fig{seasonality.png}{Seasonality curve from the Monthly Trend sheet. Each line is one year, allowing the same calendar month to be compared across years.}

\sect{Workbook map}
\begin{tabularx}{\linewidth}{@{}lX@{}}
\toprule
\textbf{Sheet} & \textbf{Purpose}\\
\midrule
Dashboard & KPI cards and charts for the executive view.\\
Monthly Trend & 48-month sales/profit table, moving average and seasonality grid.\\
Product Analysis & Sub-category performance and top/bottom product rankings.\\
Discount Analysis & Band, exact-rate, category, monthly and correlation analysis.\\
KPI Definitions & Live KPI formulas, definitions and supporting category/region tables.\\
Data & Clean order lines plus formula-derived year, quarter, ship days, returns and discount band.\\
Cleaning Log & Auditable record of every split, conversion, correction and validation check.\\
\bottomrule
\end{tabularx}

\sect{Recommendations}
\begin{enumerate}[leftmargin=1.5em,itemsep=2pt]
  \item Introduce a review or minimum-margin guardrail for discounts above 20\%,
  beginning with Tables, Bookcases and Machines.
  \item Protect peak-season demand with availability, bundles and cross-sells
  before adding broad discounts.
  \item Join returns, fulfillment cost and customer data to test whether the
  worst products and returned orders are structurally unprofitable.
\end{enumerate}

\sect{Limitations}
The discount results are associations, not causal estimates. The data covers
one retailer from 2015--2018 and does not contain cost detail, marketing
exposure or customer contact history. Return profit treatment is unknown, and
the Burlington ZIP fill is an outside factual assumption. Excel formula cells
require recalculation in Excel or LibreOffice after generation.

\notes{Added explicit validation and made the README operational rather than
placeholder-based.}{The project is reproducible, the workbook structure is
checked, and the recommendations connect directly to measured findings.}{No
causal experiment can be inferred from this observational dataset, and slicers
remain a manual PivotTable step.}{Add a controlled discount test and join
margin, cost and return outcomes before changing commercial policy.}

\sect{Deliverables}
\begin{itemize}
  \item \href{run:../superstore_beginner.xlsx}{Live Excel workbook}
  \item \href{run:../scripts/validate_project.py}{Regression validator}
  \item \href{run:beginner_report.tex}{This LaTeX source}
\end{itemize}

\end{document}
